% !TEX root = ../thesis_book_0421052099.tex
\chapter{Conclusion}\label{ch:conclusion}

\section{Summary of Contributions}\label{sec:summary}

This thesis asked what an agent gains by navigating a knowledge graph rather
than retrieving from it, and what a claim-level check against the traversed
subgraph adds on top. The answers are not uniform, and the value of the work
lies partly in where they diverge.

\Cref{sec:results} reported the measurements. AGR leads every accuracy cell of
the main matrix --- the agentic baseline hedges marginally less on CWQ, which
\Cref{sec:findings} reports as the one cell it does not lead. Every one
of the eight paired comparisons rejects, including both against a
reimplementation of the agentic baseline running on the same graph with the same
model under the same budget. That last margin is located rather than merely
claimed: on the questions where the shared cap lets the baseline finish, the
baseline is slightly ahead. The whole margin comes from the $29\%$ and
$44\%$ of questions where the cap truncates it (\Cref{sec:tog-split}). The
finding is therefore about affordability rather than about resolution quality
--- beam search cannot complete under a budget a deployed system would impose,
and guided search can. On ComplexWebQuestions AGR is the only system whose
accuracy \emph{rises} with hop count. Every other system ends below where it
started --- the strongest available form of the multi-hop claim. The agentic
baseline is the one that does not simply decay: it falls and then partially
recovers. That recovery is weaker than a rise and is reported as such in
\Cref{sec:hop-results}. AGR runs at half the language-model calls of the agentic
baseline and never exhausts the call budget that clips that baseline on $44\%$
of multi-hop questions. Structural hallucination is eliminated: over both
datasets, zero ungrounded assertions in $1{,}709$ asserted entities, against
$22.1\%$ of the $1{,}001$ the parametric control asserts. And three of four
ablations detect nothing. That null result is reported at the same length as
the one that does.

\Cref{sec:erroranalysis} read all $259$ remaining failures and named each
mechanism. The findings that generalise beyond this system are three.
\emph{Context-stripping}: decomposing a question can discard a qualifier that
appears only in a later clause, so an early sub-objective resolves the wrong
sense of an ambiguous term before the disambiguating information arrives --- and
in the strongest specimen, decomposition actively suppressed the one anchor that
could have discriminated, driving its score to the noise floor. \emph{The echo
attractor}: the most common failure of a graph-navigating system is not to
invent an entity but to surface a real, well-grounded, verifiably-true entity
that sits one node away from the answer. The verifier passes it because it
genuinely is grounded. And \emph{shared attractors break consensus-based
evaluation}: five independent systems converge on the same wrong answer often
enough that a policy of rescoring on majority agreement would have converted
this failure mode into apparent correctness.

Four contributions follow from that body of work. \Cref{sec:contribution}
claimed six. The other two are stated above rather than repeated as items here
--- the echo attractor is the second of the three findings just given, and the
pre-registered evaluation protocol is reported in the \Cref{sec:setup}
paragraph, including the threshold it missed.

\textbf{The framework and its verification layer.} An agentic KGQA system that
returns its answer paired with the triples supporting it, built so that the
supporting evidence is a byproduct of the control flow rather than a post-hoc
reconstruction.

\textbf{A component-level ablation.} Four single-deviation conditions with
paired significance testing, closing the attribution gap of
\Cref{sec:positioning} for this architecture --- and finding that only one of
the four components has a detectable accuracy effect.

\textbf{Stratum-dependent decomposition.} The literature treats decomposition as
straightforwardly beneficial. It is not: removing the planner \emph{improves}
WebQSP accuracy by $0.083$ F1 at $p = 0.006$ while cutting tokens by $31\%$, and
trends toward harming CWQ. The asymmetry, not the pooled average, is the result.
It says that decomposition should be gated on question structure rather than
applied unconditionally.

\textbf{Quantified benchmark defect rates.} Fifty-seven questions across the two
test samples where the benchmark, not the system, needed correcting --- with the
two evidence signatures that distinguish a world-fact error from an
annotation-pipeline error, and the provenance argument that explains where one
class of them comes from.

\section{Limitations}\label{sec:limitations-final}

\Cref{sec:threats} states the threats to validity in full, against the
measurements they bear on. Four of them constrain what may be concluded here
more than the rest, and are repeated for that reason.

\textbf{Wrongful acceptance is unmeasured.} The verifier logs the claims it
rejects and not the ones it accepts, so one polarity of its error is reported at
census scale and the other from a single manually reconstructed specimen. The
exposed population is bounded only as $[39, 2{,}008]$ of $2{,}008$ accepted
claims (\Cref{sec:verify-failure-modes}), and nothing measured here narrows it.

\textbf{Structural grounding is not answer adequacy.} Both navigating systems
assert zero ungrounded entities on WebQSP, which says the answers are made of
the graph's own material; it does not say they answer the question. The echo
attractor (\Cref{sec:echo}) is the counter-example, and it is the failure this
census finds most often.

\textbf{The evaluation is single-environment, single-backbone, and single-run.}
Absolute levels are bounded by an environment built from answer-containing
subgraphs, and run-to-run variance is unmeasured against a backbone that repeats
a trajectory on about two probes in three.

\textbf{The baseline comparisons are bounded in the baselines' favour, and
against them.} The agentic baseline prunes from a narrower candidate set than
AGR and the static one retrieves a single logical hop with an unordered fanout
cap, so the reported margins are lower bounds on what equal-width, equal-radius
reimplementations would show (\Cref{sec:future-baselines}). Entity linking is
assumed throughout, topic entities being given by the benchmarks.

\section{Future Work}\label{sec:future}

The census makes it possible to separate proposals backed by a counted, repeated
pattern from proposals backed by an argument. Both are listed. They are not the
same kind of claim.

\subsection{Repairs the Census Earned}

Three defects each recur with a uniform mechanism across unrelated questions,
and each is small. \emph{Entity extraction} (nine instances): the step
converting drafted answer text into the scored entity list defaults to the
sentence's grammatical subject rather than its predicate values, and fixing it
requires telling the claim decomposer which argument of the drafted sentence is
the answer --- an instruction, not an architecture. \emph{Evaluator abandonment}
(five instances): relations scoring as high as $0.73$ are explored, backtracked
away from, and never revisited, across four unrelated domains, which reads as a
threshold or backtracking-policy defect and would be diagnosed by logging the
evaluator's decision against the score of the relation it discards.
\emph{Drafting commitment} (six instances, all six of the \textsf{other}
category): when the evaluator has resolved the exact gold value and the verifier
has returned \textsf{grounded}, the drafter should commit to it rather than
stating the fact and declaring it undeterminable in the same sentence.

\subsection{Logging Accepted Claims to Make Wrongful Acceptance
Measurable}\label{sec:future-logging}

Persist each accepted claim together with the triples that matched it. This
converts\index{verification layer} the wrongly-accepted class from an anecdote
into a rate with a denominator, and makes grounded-but-wrong answers
self-diagnosing from the logs rather than by manual trace reconstruction. The
same change closes a second gap: the run record currently stores the
\emph{number} of supporting triples rather than the triples
(\Cref{sec:output-contract}), so the output contract holds inside the run but
cannot be checked afterwards from any committed file. It is the highest-value
change in this list relative to its cost, and it is a logging change.

\subsection{Operators the Failure Categories Specify}

Three additions are specified by categories rather than motivated by them.
\emph{Adaptive decomposition gating}: a decomposition step should be gated on
whether the rewritten sub-objective preserves the question's actual
interrogative --- the Phoenician case of \Cref{sec:decomposition-errors} lost
the word ``where'' entirely --- on whether the step resolves an entity that was
already unambiguous, and on whether the sub-objectives reference each other at
all, since ones that do not are descriptors of a single entity rather than a
chain. \emph{A set-intersection operator}: \textsf{composite\_claim} is $47$
cases and CWQ's largest category at $46$ of them, and its
\textsf{no\_set\_intersection} subtype names the gap precisely --- the
architecture has no AND-join primitive, so a ``find the overlap'' sub-objective
degrades into another relation-scoring pass over an incomplete first-hop set.
\emph{Ordinal and literal-value support}: \textsf{kg\_gap} is $44$ cases and
CWQ's largest hedge category at $24$, and two of its subtypes need a
minimum-or-maximum-by-value operation and a literal-extraction path rather than
a bigger graph.

\subsection{Semantic-Level Verification}

All five wrongly-rejected cases share one mechanism: the structural check
requires a single edge asserting the drafted claim, and a transitively-true
multi-hop fact does not present one. A check that accepts a bounded relation
chain expressing the claim would recover those answers. The relation-path
formulation of RoG~\cite{luo2024reasoning} is the natural source for what such a
chain may look like, and using it as an acceptance criterion rather than as a
retrieval plan is the change proposed. The Tier-2 judge of
\Cref{sec:groundedness-metric} is a prototype of exactly this judgement, which
suggests folding a bounded form of it into the loop; its $\kappa$ of $0.6995$ is
also a caution about how reliable such a check would be.

\subsection{Measurements This Thesis Could Not
Make}\label{sec:future-baselines}

\label{sec:future-path-fidelity}%
\emph{Path fidelity} is the third evaluation criterion of the approved proposal,
deferred for the reason given in \Cref{sec:metrics-not-run}: it needs the
benchmarks' gold SPARQL relation chains, which the distribution used here does
not carry. Reconstructing them would measure whether the system reaches the
right answer \emph{by the right path} --- a question the groundedness metrics
can only approximate, and one that bears directly on the echo attractor.

Two baseline bounds are cheap to lift, and both are re-runs rather than
redesigns. \emph{Equalise the candidate widths}: re-run the agentic baseline at
AGR's $300$ relations and $200$ neighbours, leaving beam width, depth, budget,
prompts, and scoring untouched. That decides whether the unclipped-subset
comparison of \Cref{sec:tog-split} measured search strategy or candidate supply,
and both outcomes are usable --- though it will also raise the clip rate,
because wider candidate sets make each pruning call dearer, so the split must be
reported rather than the aggregate. \emph{Widen the static retriever's radius,
and order its fanout}: re-run it with neighbours expanded a second time,
everything else frozen, and report the same strata, so that its collapse on deep
CWQ questions can be attributed to the paradigm rather than to the particular
radius chosen. The same pass should impose an ordering on the fanout cap, so
that the $100$ edges kept per topic entity are a stated selection rather than an
arbitrary one (\Cref{sec:baseline-graphrag}). Until then those two limits are
confounded and only their sum is measured.

\subsection{Beyond This Environment}

Four directions extend the work rather than repairing it. \emph{LLM-constructed
knowledge graphs} as a controlled comparison would separate what the framework
contributes from what a curated environment contributes, by holding the agent
fixed while varying the substrate's provenance and noise profile.
\emph{Rollout-based value estimation} would supply what
\Cref{sec:mcts-clarification} explicitly disclaims: the current scorer evaluates
each candidate once and never revises it in light of what is found later, and
genuine tree search over this same tool API is a well-defined next system rather
than a rebranding of this one. \emph{Temporal and multi-source environments}
address the failure family that recurs throughout \Cref{sec:erroranalysis} ---
questions with a time qualifier the graph cannot express --- and raise the
question of what a claim check means when sources disagree. And \emph{transfer
to high-stakes domains} is where the auditability property, rather than the
accuracy, is the point: an answer paired with the triples supporting it is worth
more in a setting where a reader must be able to check it than in one where the
reader only wants the answer.
\section{Closing Remark}

The result this work would most want carried forward is not the accuracy margin.
It is the shape of the gap that remains. A system that traverses a real graph
and copies its answers out of real triples can be made to assert nothing
ungrounded --- that problem is solved, and it is solved by the navigation
paradigm rather than by any checking layer built on top of it. What remains, and
what this thesis's census found more often than any other mechanism, is an
answer that is grounded, connected, and true, and is not the answer to the
question that was asked. Fact verification against a knowledge graph turns out
to be the easier half of the problem; \emph{relevance} verification is the half
still open.

\endinput
